%! Unidad 5 — Nomenclatura de química inorgánica: ejercicios resueltos y ejercicios finales
%! Fuente: Fuentes/Quimica/Unidad 5 («Ejercicios», «Ejercicios resueltos» y «Practica y calcula»).
%! Enunciados adaptados con compuestos cambiados; soluciones rehechas. Dificultad: \(\star\) baja, \(\star\star\) media, \(\star\star\star\) alta.

\sloppy % nombres largos sin guiones: se permite más espacio entre palabras (\fussy al final)

\section{Ejercicios resueltos}\label{sec:ud5-resueltos}

\begin{example}
    \textbf{Identificar el tipo de compuesto.} Clasifica estas sustancias y nómbralas: \ch{Ba(OH)2}, \ch{N2O4}, \ch{HIO4},
    \ch{K2O2}, \ch{NH4Cl}, \ch{SbH3}.

    \textbf{Solución.} Para clasificar, mira primero cuántos elementos hay y qué hay a la derecha de la fórmula:
    \begin{itemize}
        \item \ch{Ba(OH)2}: metal y \ch{OH}: \textbf{hidróxido}. Hidróxido de bario (el bario solo tiene \(+2\)).
        \item \ch{N2O4}: binario con O a la derecha: \textbf{óxido}. Tetraóxido de dinitrógeno; óxido de nitrógeno(IV)
        (\(2x - 8 = 0\), \(x = +4\)). No se simplifica la fórmula: el \ch{N2O4} y el \ch{NO2} son sustancias distintas.
        \item \ch{HIO4}: H, un no metal y O: \textbf{oxoácido}. Ácido peryódico (I con \(+7\)); hidrogeno(tetraoxidoyodato);
        \ch{IO3(OH)}, hidroxidotrioxidoyodo.
        \item \ch{K2O2}: el K solo tiene \(+1\), así que cada O tiene \(-1\): \textbf{peróxido}. Peróxido de potasio; dióxido de
        dipotasio.
        \item \ch{NH4Cl}: catión amonio y anión cloruro: \textbf{sal}, aunque sea ternaria se nombra como una binaria. Cloruro de
        amonio.
        \item \ch{SbH3}: binario con H a la derecha: \textbf{hidruro}. Trihidruro de antimonio; hidruro de antimonio(III);
        estibano.
    \end{itemize}
\end{example}

\begin{example}
    \textbf{Formular a partir del número de oxidación.} Formula el óxido de cloro(VII), el sulfuro de plomo(IV), el hidróxido de
    hierro(II) y el peróxido de aluminio.

    \textbf{Solución.} En cada caso escribimos los iones (o los estados de oxidación) y aplicamos la regla de la suma:
    \begin{itemize}
        \item Óxido de cloro(VII): ¡cuidado!, el O aparece antes que el Cl en la secuencia, así que este compuesto es un haluro de
        oxígeno. Con Cl \(+7\) y O \(-2\) la proporción es 2:7: la fórmula es \ch{O7Cl2}, dicloruro de heptaoxígeno. El nombre
        «óxido de cloro(VII)» no es correcto.
        \item Sulfuro de plomo(IV): \ch{Pb^4+} y \ch{S^2-}: \ch{Pb2S4} \(\rightarrow\) \ch{PbS2}.
        \item Hidróxido de hierro(II): \ch{Fe^2+} y \ch{OH-}: \ch{Fe(OH)2}.
        \item Peróxido de aluminio: \ch{Al^3+} y \ch{O2^2-}. Con 2 Al y 3 \ch{O2^2-}: \ch{Al2(O2)3} \(=\) \ch{Al2O6}. Como es un
        peróxido, no se simplifica.
    \end{itemize}
\end{example}

\begin{example}
    \textbf{Nombrar un oxoácido «desconocido».} Nombra el \ch{H2MnO4} en las nomenclaturas de hidrógeno y de adición. ¿Qué
    estado de oxidación tiene el manganeso? ¿Cómo se llama el anión que se forma al perder sus dos \ch{H+}?

    \textbf{Solución.} \textbf{1. Hidrógeno}: 2 H, 4 O, 1 Mn: dihidrogeno(tetraoxidomanganato). \textbf{2. Adición}:
    \ch{MnO2(OH)2}, dihidroxidodioxidomanganeso. \textbf{3. Estado de oxidación}: \(2 + x - 8 = 0\), \(x = +6\).
    \textbf{4. Anión}: \ch{MnO4^2-}, tetraoxidomanganato(2\nobreakdash--), cuyo nombre tradicional es \textbf{manganato}. No lo confundas con
    el \ch{MnO4-}, permanganato (Mn \(+7\)).
\end{example}

\begin{example}
    \textbf{Nombrar una oxosal sin saber la carga del catión.} Nombra el \ch{Co(ClO2)3} y el \ch{NiSO3} en los tres sistemas.

    \textbf{Solución.} \textbf{\ch{Co(ClO2)3}}: hay 3 aniones \ch{ClO2}. Si su carga fuera \(-1\), el Cl tendría \(x - 4 = -1\),
    \(x = +3\), correcto, y el catión sería \ch{Co^3+}, también posible. Es el clorito de cobalto(III);
    tris(dioxidoclorato) de cobalto; dioxidoclorato(1\nobreakdash--) de cobalto(3+).

    \textbf{\ch{NiSO3}}: un catión y un anión. Si el anión fuera \ch{SO3^2-} (S con \(+4\)), el catión sería \ch{Ni^2+}, que es
    posible. Si fuera \ch{SO3^3-}, el S tendría \(+3\), que no es válido. Es el sulfito de níquel(II); trioxidosulfato de níquel;
    trioxidosulfato(2\nobreakdash--) de níquel(2+).
\end{example}

\begin{example}
    \textbf{Sales ácidas e hidratadas.} a) Formula el hidrogenofosfato de amonio y el bis[dihidrogeno(tetraoxidofosfato)] de
    magnesio. b) Nombra el \ch{Na2B4O7 * 10 H2O} (bórax), sabiendo que su anión es el heptaoxidotetraborato(2\nobreakdash--).

    \textbf{Solución.} a) Hidrogenofosfato: \ch{HPO4^2-}; amonio: \ch{NH4+}. Hacen falta dos amonios: \ch{(NH4)2HPO4}. Es un
    abono muy usado. El segundo tiene dos aniones \ch{H2PO4-} y un magnesio: \ch{Mg(H2PO4)2}.

    b) Heptaoxidotetraborato de disodio---agua (1/10), o heptaoxidotetraborato de disodio decahidrato.
\end{example}

%========================================================================================
\section{Ejercicios de aplicación}\label{sec:ud5-aplicacion}
%========================================================================================

\subsection{Conceptos previos}

\begin{exercise}
    \(\star\) Clasifica estas sustancias en simples, binarias, ternarias o cuaternarias: \ch{Cl2}, \ch{KNO3}, \ch{CaO}, \ch{NaHSO4},
    \ch{Fe}, \ch{H3PO4}, \ch{HBr}, \ch{NH4HCO3}, \ch{O3}, \ch{Mg(OH)2}.
    \begin{solu}
        \textcolor{red}{Simples}: \ch{Cl2}, \ch{Fe}, \ch{O3}. \textcolor{red}{Binarias}: \ch{CaO}, \ch{HBr}.
        \textcolor{red}{Ternarias}: \ch{KNO3}, \ch{H3PO4}, \ch{Mg(OH)2}. \textcolor{red}{Cuaternarias}: \ch{NaHSO4}, \ch{NH4HCO3}.
    \end{solu}
\end{exercise}

\begin{exercise}
    \(\star\star\) Calcula el estado de oxidación de cada elemento en: a) \ch{HNO2}; b) \ch{Na2CrO4}; c) \ch{Al2(SO4)3}; d)
    \ch{MnO4-}; e) \ch{H2O2}; f) \ch{CaH2}; g) \ch{NH3}; h) \ch{S2O3^2-}.
    \begin{solu}
        a) \textcolor{red}{H \(+1\), N \(+3\), O \(-2\)}. b) \textcolor{red}{Na \(+1\), Cr \(+6\), O \(-2\)}. c)
        \textcolor{red}{Al \(+3\), S \(+6\), O \(-2\)} (anión \ch{SO4^2-}). d) \textcolor{red}{Mn \(+7\), O \(-2\)}. e)
        \textcolor{red}{H \(+1\), O \(-1\)}. f) \textcolor{red}{Ca \(+2\), H \(-1\)}. g) \textcolor{red}{N \(-3\), H \(+1\)}.
        h) \(2x - 6 = -2\): \textcolor{red}{S \(+2\), O \(-2\)}.
    \end{solu}
\end{exercise}

\begin{exercise}
    \(\star\star\) Con la regla de la suma, decide si son posibles estos compuestos y corrige los que no lo sean: a) \ch{NaS};
    b) \ch{BaCl2}; c) \ch{AlO}; d) \ch{K3P}; e) \ch{MgN}; f) \ch{ZnBr}.
    \begin{solu}
        a) \textcolor{red}{No}: \(+1 - 2 = -1\); \textcolor{red}{\ch{Na2S}}. b) \textcolor{red}{Sí}. c) \textcolor{red}{No}:
        \(+3 - 2 \neq 0\); \textcolor{red}{\ch{Al2O3}}. d) \textcolor{red}{Sí}: \(3 - 3 = 0\). e) \textcolor{red}{No}:
        \(+2 - 3 \neq 0\); \textcolor{red}{\ch{Mg3N2}}. f) \textcolor{red}{No}; \textcolor{red}{\ch{ZnBr2}}.
    \end{solu}
\end{exercise}

\begin{exercise}
    \(\star\star\) Usando la secuencia de los elementos (figura~\ref{fig:ud5-secuencia}), escribe correctamente la fórmula de los
    compuestos formados por: a) cloro y oxígeno, con 2 Cl y 1 O; b) flúor y yodo, con 5 F; c) nitrógeno e hidrógeno; d) azufre y
    flúor, con 6 F; e) silicio y carbono; f) bromo y cloro, uno de cada.
    \begin{solu}
        a) \textcolor{red}{\ch{OCl2}}; b) \textcolor{red}{\ch{IF5}}; c) \textcolor{red}{\ch{NH3}}; d) \textcolor{red}{\ch{SF6}};
        e) \textcolor{red}{\ch{SiC}}; f) \textcolor{red}{\ch{BrCl}} (el Br aparece antes que el Cl, porque dentro del grupo se avanza
        de abajo arriba).
    \end{solu}
\end{exercise}

\subsection{Sustancias simples e iones}

\begin{exercise}
    \(\star\) Nombra o formula: a) \ch{Kr}; b) \ch{Cl2}; c) \ch{Pt}; d) \ch{P4}; e) \ch{H2}; f) xenón; g) plomo; h) trioxígeno;
    i) dinitrógeno; j) octaazufre.
    \begin{solu}
        a) \textcolor{red}{criptón}; b) \textcolor{red}{dicloro}; c) \textcolor{red}{platino}; d) \textcolor{red}{tetrafósforo};
        e) \textcolor{red}{dihidrógeno}; f) \textcolor{red}{\ch{Xe}}; g) \textcolor{red}{\ch{Pb}}; h) \textcolor{red}{\ch{O3}};
        i) \textcolor{red}{\ch{N2}}; j) \textcolor{red}{\ch{S8}}.
    \end{solu}
\end{exercise}

\begin{exercise}
    \(\star\) Nombra o formula estos iones: a) \ch{Li+}; b) \ch{S^2-}; c) \ch{Ni^2+}; d) \ch{CN-}; e) \ch{Te^2-}; f) \ch{Hg2^2+};
    g) cloruro; h) aluminio(3+); i) nitruro; j) oxonio; k) cromo(3+); l) peróxido.
    \begin{solu}
        a) \textcolor{red}{litio(1+)}; b) \textcolor{red}{sulfuro}; c) \textcolor{red}{níquel(2+)}; d) \textcolor{red}{cianuro};
        e) \textcolor{red}{telururo}; f) \textcolor{red}{dimercurio(2+)}; g) \textcolor{red}{\ch{Cl-}}; h)
        \textcolor{red}{\ch{Al^3+}}; i) \textcolor{red}{\ch{N^3-}}; j) \textcolor{red}{\ch{H3O+}}; k) \textcolor{red}{\ch{Cr^3+}};
        l) \textcolor{red}{\ch{O2^2-}}.
    \end{solu}
\end{exercise}

\subsection{Compuestos binarios con hidrógeno}

\begin{exercise}
    \(\star\) Nombra de todas las maneras posibles: a) \ch{RbH}; b) \ch{SrH2}; c) \ch{NiH2}; d) \ch{PbH2}; e) \ch{HI}; f)
    \ch{H2Te}; g) \ch{BiH3}; h) \ch{GeH4}; i) \ch{AlH3}; j) \ch{MnH2}.
    \begin{solu}
        a) \textcolor{red}{hidruro de rubidio}. b) \textcolor{red}{dihidruro de estroncio; hidruro de estroncio}. c)
        \textcolor{red}{dihidruro de níquel; hidruro de níquel(II); hidruro de níquel(2+)}. d) \textcolor{red}{dihidruro de plomo;
        hidruro de plomo(II); hidruro de plomo(2+)}. e) \textcolor{red}{yoduro de hidrógeno; yodano}. f) \textcolor{red}{telururo
        de dihidrógeno; telururo de hidrógeno; telano}. g) \textcolor{red}{trihidruro de bismuto; hidruro de bismuto(III);
        bismutano}. h) \textcolor{red}{tetrahidruro de germanio; hidruro de germanio(IV); germano}. i) \textcolor{red}{trihidruro
        de aluminio; hidruro de aluminio; alumano}. j) \textcolor{red}{dihidruro de manganeso; hidruro de manganeso(II); hidruro
        de manganeso(2+)}.
    \end{solu}
\end{exercise}

\begin{exercise}
    \(\star\) Formula: a) hidruro de cesio; b) hidruro de oro(III); c) hidruro de cobalto(3+); d) dihidruro de calcio; e) bromuro
    de hidrógeno; f) seleniuro de hidrógeno; g) fosfano; h) arsano; i) silano; j) hidruro de cromo(II); k) borano; l) sulfano.
    \begin{solu}
        a) \textcolor{red}{\ch{CsH}}; b) \textcolor{red}{\ch{AuH3}}; c) \textcolor{red}{\ch{CoH3}}; d) \textcolor{red}{\ch{CaH2}};
        e) \textcolor{red}{\ch{HBr}}; f) \textcolor{red}{\ch{H2Se}}; g) \textcolor{red}{\ch{PH3}}; h) \textcolor{red}{\ch{AsH3}};
        i) \textcolor{red}{\ch{SiH4}}; j) \textcolor{red}{\ch{CrH2}}; k) \textcolor{red}{\ch{BH3}}; l) \textcolor{red}{\ch{H2S}}.
    \end{solu}
\end{exercise}

\begin{exercise}
    \(\star\star\) ¿Por qué el \ch{HCl} se escribe con el hidrógeno a la izquierda y el \ch{NaH} con el hidrógeno a la derecha?
    ¿Qué estado de oxidación tiene el hidrógeno en cada caso? ¿Por qué no es correcto llamar «ácido clorhídrico» al \ch{HCl} puro?
    \begin{solu}
        \textcolor{red}{En la secuencia, el H aparece antes que el Cl y después que el Na}: en el \ch{HCl} hace de electropositivo
        (\textcolor{red}{\(+1\)}) y en el \ch{NaH}, de electronegativo (\textcolor{red}{\(-1\)}). \textcolor{red}{«Ácido
        clorhídrico» es el nombre de su disolución acuosa}, no del compuesto: el \ch{HCl} puro es un gas, el cloruro de hidrógeno.
    \end{solu}
\end{exercise}

\subsection{Compuestos binarios con oxígeno}

\begin{exercise}
    \(\star\) Nombra de todas las maneras posibles: a) \ch{K2O}; b) \ch{CaO}; c) \ch{FeO}; d) \ch{CuO}; e) \ch{Cr2O3}; f)
    \ch{P2O3}; g) \ch{SeO2}; h) \ch{I2O5}; i) \ch{Ag2O}; j) \ch{TiO2}; k) \ch{Mn2O7}; l) \ch{NO}.
    \begin{solu}
        a) \textcolor{red}{óxido de dipotasio; óxido de potasio}. b) \textcolor{red}{óxido de calcio}. c) \textcolor{red}{monóxido
        de hierro; óxido de hierro(II); óxido de hierro(2+)}. d) \textcolor{red}{monóxido de cobre; óxido de cobre(II); óxido de
        cobre(2+)}. e) \textcolor{red}{trióxido de dicromo; óxido de cromo(III); óxido de cromo(3+)}. f) \textcolor{red}{trióxido de
        difósforo; óxido de fósforo(III)}. g) \textcolor{red}{dióxido de selenio; óxido de selenio(IV)}. h) \textcolor{red}{Debe
        escribirse \ch{O5I2}: diyoduro de pentaoxígeno}. i) \textcolor{red}{óxido de diplata; óxido de plata}. j)
        \textcolor{red}{dióxido de titanio; óxido de titanio(IV); óxido de titanio(4+)}. k) \textcolor{red}{heptaóxido de
        dimanganeso; óxido de manganeso(VII); óxido de manganeso(7+)}. l) \textcolor{red}{monóxido de nitrógeno; óxido de
        nitrógeno(II)}.
    \end{solu}
\end{exercise}

\begin{exercise}
    \(\star\) Formula: a) óxido de bario; b) óxido de níquel(III); c) dióxido de nitrógeno; d) óxido de cobre(1+); e) trióxido de
    azufre; f) óxido de arsénico(V); g) óxido de estaño(II); h) monóxido de carbono; i) óxido de aluminio; j) dibromuro de
    oxígeno; k) óxido de plomo(II); l) óxido de cobalto(III).
    \begin{solu}
        a) \textcolor{red}{\ch{BaO}}; b) \textcolor{red}{\ch{Ni2O3}}; c) \textcolor{red}{\ch{NO2}}; d) \textcolor{red}{\ch{Cu2O}};
        e) \textcolor{red}{\ch{SO3}}; f) \textcolor{red}{\ch{As2O5}}; g) \textcolor{red}{\ch{SnO}}; h) \textcolor{red}{\ch{CO}};
        i) \textcolor{red}{\ch{Al2O3}}; j) \textcolor{red}{\ch{OBr2}}; k) \textcolor{red}{\ch{PbO}}; l) \textcolor{red}{\ch{Co2O3}}.
    \end{solu}
\end{exercise}

\begin{exercise}
    \(\star\star\) Nombra de todas las maneras posibles estos peróxidos: a) \ch{Rb2O2}; b) \ch{MgO2}; c) \ch{ZnO2}; d) \ch{BaO2};
    e) \ch{Ag2O2}; f) \ch{Ni2O6}; g) \ch{HgO2}.
    \begin{solu}
        a) \textcolor{red}{dióxido de dirrubidio; peróxido de rubidio; dióxido(2\nobreakdash--) de rubidio}. b) \textcolor{red}{dióxido de
        magnesio; peróxido de magnesio}. c) \textcolor{red}{dióxido de cinc; peróxido de cinc}. d) \textcolor{red}{dióxido de bario;
        peróxido de bario}. e) \textcolor{red}{dióxido de diplata; peróxido de plata}. f) \textcolor{red}{hexaóxido de diníquel;
        peróxido de níquel(III); dióxido(2\nobreakdash--) de níquel(3+)}. g) \textcolor{red}{dióxido de mercurio; peróxido de mercurio(II);
        dióxido(2\nobreakdash--) de mercurio(2+)}.
    \end{solu}
\end{exercise}

\begin{exercise}
    \(\star\star\) Formula: a) peróxido de litio; b) dióxido de disodio; c) peróxido de cobre(I); d) peróxido de estroncio; e)
    peróxido de hierro(III); f) dióxido(2\nobreakdash--) de cadmio; g) peróxido de cromo(III); h) agua oxigenada.
    \begin{solu}
        a) \textcolor{red}{\ch{Li2O2}}; b) \textcolor{red}{\ch{Na2O2}}; c) \textcolor{red}{\ch{Cu2O2}}; d) \textcolor{red}{\ch{SrO2}};
        e) \textcolor{red}{\ch{Fe2O6}}; f) \textcolor{red}{\ch{CdO2}}; g) \textcolor{red}{\ch{Cr2O6}}; h) \textcolor{red}{\ch{H2O2}}.
    \end{solu}
\end{exercise}

\begin{exercise}
    \(\star\star\) El \ch{PbO2} y el \ch{BaO2} tienen fórmulas parecidas. ¿Son los dos peróxidos? Razónalo con los estados de
    oxidación y nómbralos.
    \begin{solu}
        \textcolor{red}{No}. El bario solo tiene \(+2\): para que la suma sea cero, cada O tiene \(-1\) y el \ch{BaO2} es
        \textcolor{red}{peróxido de bario}. El plomo puede tener \(+4\), con O en \(-2\): el \ch{PbO2} es un óxido,
        \textcolor{red}{dióxido de plomo u óxido de plomo(IV)}.
    \end{solu}
\end{exercise}

\subsection{Sales binarias}

\begin{exercise}
    \(\star\) Nombra de todas las maneras posibles: a) \ch{NaI}; b) \ch{MgCl2}; c) \ch{FeS}; d) \ch{SnBr2}; e) \ch{AuCl3}; f)
    \ch{CuI2}; g) \ch{Ag2S}; h) \ch{CaF2}; i) \ch{PtCl4}; j) \ch{Hg2Cl2}; k) \ch{AlN}; l) \ch{K2Se}.
    \begin{solu}
        a) \textcolor{red}{yoduro de sodio}. b) \textcolor{red}{dicloruro de magnesio; cloruro de magnesio}. c)
        \textcolor{red}{sulfuro de hierro; sulfuro de hierro(II); sulfuro de hierro(2+)}. d) \textcolor{red}{dibromuro de estaño;
        bromuro de estaño(II); bromuro de estaño(2+)}. e) \textcolor{red}{tricloruro de oro; cloruro de oro(III); cloruro de
        oro(3+)}. f) \textcolor{red}{diyoduro de cobre; yoduro de cobre(II); yoduro de cobre(2+)}. g) \textcolor{red}{sulfuro de
        diplata; sulfuro de plata}. h) \textcolor{red}{difluoruro de calcio; fluoruro de calcio}. i) \textcolor{red}{tetracloruro
        de platino; cloruro de platino(IV); cloruro de platino(4+)}. j) \textcolor{red}{dicloruro de dimercurio; cloruro de
        mercurio(I)} (contiene el ion \ch{Hg2^2+}). k) \textcolor{red}{nitruro de aluminio}. l) \textcolor{red}{seleniuro de
        dipotasio; seleniuro de potasio}.
    \end{solu}
\end{exercise}

\begin{exercise}
    \(\star\) Formula: a) bromuro de litio; b) cloruro de cobalto(II); c) sulfuro de cinc; d) yoduro de plomo(II); e) fluoruro de
    aluminio; f) sulfuro de níquel(III); g) cloruro de cromo(3+); h) nitruro de sodio; i) tetrabromuro de estaño; j) fosfuro de
    calcio; k) yoduro de mercurio(II); l) carburo de aluminio (con \ch{C^4-}).
    \begin{solu}
        a) \textcolor{red}{\ch{LiBr}}; b) \textcolor{red}{\ch{CoCl2}}; c) \textcolor{red}{\ch{ZnS}}; d) \textcolor{red}{\ch{PbI2}};
        e) \textcolor{red}{\ch{AlF3}}; f) \textcolor{red}{\ch{Ni2S3}}; g) \textcolor{red}{\ch{CrCl3}}; h) \textcolor{red}{\ch{Na3N}};
        i) \textcolor{red}{\ch{SnBr4}}; j) \textcolor{red}{\ch{Ca3P2}}; k) \textcolor{red}{\ch{HgI2}}; l) \textcolor{red}{\ch{Al4C3}}.
    \end{solu}
\end{exercise}

\begin{exercise}
    \(\star\star\) Nombra o formula estos compuestos entre no metales: a) \ch{PBr3}; b) \ch{CCl4}; c) \ch{SiF4}; d) \ch{ClF3};
    e) \ch{As2S3}; f) \ch{BCl3}; g) pentacloruro de fósforo; h) disulfuro de carbono; i) cloruro de yodo(I); j) hexafluoruro de
    selenio; k) nitruro de boro; l) tribromuro de arsénico.
    \begin{solu}
        a) \textcolor{red}{tribromuro de fósforo; bromuro de fósforo(III)}. b) \textcolor{red}{tetracloruro de carbono; cloruro de
        carbono(IV)}. c) \textcolor{red}{tetrafluoruro de silicio; fluoruro de silicio}. d) \textcolor{red}{trifluoruro de cloro;
        fluoruro de cloro(III)}. e) \textcolor{red}{trisulfuro de diarsénico; sulfuro de arsénico(III)}. f)
        \textcolor{red}{tricloruro de boro; cloruro de boro}. g) \textcolor{red}{\ch{PCl5}}; h) \textcolor{red}{\ch{CS2}}; i)
        \textcolor{red}{\ch{ICl}}; j) \textcolor{red}{\ch{SeF6}}; k) \textcolor{red}{\ch{BN}}; l) \textcolor{red}{\ch{AsBr3}}.
    \end{solu}
\end{exercise}

\begin{exercise}
    \(\star\star\) Un alumno ha escrito estos nombres. Corrígelos y justifica la corrección: a) \ch{NaCl}, «cloruro de sodio(I)»;
    b) \ch{PCl3}, «cloruro de fósforo(3+)»; c) \ch{FeCl3}, «cloruro de hierro»; d) \ch{Cl2O3}, «trióxido de dicloro».
    \begin{solu}
        a) \textcolor{red}{Cloruro de sodio}: el sodio solo tiene un estado de oxidación y no hace falta indicarlo (aunque no es
        incorrecto). b) \textcolor{red}{Cloruro de fósforo(III)} o tricloruro de fósforo: no es iónico y no admite número de carga.
        c) \textcolor{red}{Cloruro de hierro(III)} o tricloruro de hierro: el hierro tiene dos estados de oxidación y hay que
        indicarlo. d) \textcolor{red}{\ch{O3Cl2}, dicloruro de trioxígeno}: el O va antes que el Cl en la secuencia, así que se
        escribe a la izquierda y el compuesto es un haluro de oxígeno, no un óxido.
    \end{solu}
\end{exercise}

\subsection{Hidróxidos}

\begin{exercise}
    \(\star\) Nombra de todas las maneras posibles: a) \ch{CsOH}; b) \ch{Sr(OH)2}; c) \ch{Fe(OH)2}; d) \ch{Cu(OH)2}; e)
    \ch{Pb(OH)4}; f) \ch{Cd(OH)2}; g) \ch{Mn(OH)3}; h) \ch{NH4OH}.
    \begin{solu}
        a) \textcolor{red}{hidróxido de cesio}. b) \textcolor{red}{dihidróxido de estroncio; hidróxido de estroncio}. c)
        \textcolor{red}{dihidróxido de hierro; hidróxido de hierro(II); hidróxido de hierro(2+)}. d) \textcolor{red}{dihidróxido de
        cobre; hidróxido de cobre(II); hidróxido de cobre(2+)}. e) \textcolor{red}{tetrahidróxido de plomo; hidróxido de plomo(IV);
        hidróxido de plomo(4+)}. f) \textcolor{red}{dihidróxido de cadmio; hidróxido de cadmio}. g) \textcolor{red}{trihidróxido
        de manganeso; hidróxido de manganeso(III); hidróxido de manganeso(3+)}. h) \textcolor{red}{hidróxido de amonio}.
    \end{solu}
\end{exercise}

\begin{exercise}
    \(\star\) Formula: a) hidróxido de potasio; b) hidróxido de aluminio; c) hidróxido de níquel(III); d) hidróxido de oro(I);
    e) dihidróxido de berilio; f) hidróxido de estaño(4+); g) hidróxido de cromo(II); h) trihidróxido de bismuto.
    \begin{solu}
        a) \textcolor{red}{\ch{KOH}}; b) \textcolor{red}{\ch{Al(OH)3}}; c) \textcolor{red}{\ch{Ni(OH)3}}; d)
        \textcolor{red}{\ch{AuOH}}; e) \textcolor{red}{\ch{Be(OH)2}}; f) \textcolor{red}{\ch{Sn(OH)4}}; g)
        \textcolor{red}{\ch{Cr(OH)2}}; h) \textcolor{red}{\ch{Bi(OH)3}}.
    \end{solu}
\end{exercise}

\subsection{Oxoácidos}

\begin{exercise}
    \(\star\) Nombra con la nomenclatura de hidrógeno: a) \ch{H2CO3}; b) \ch{HClO3}; c) \ch{HMnO4}; d) \ch{HBrO}; e) \ch{H2SO4};
    f) \ch{H4SiO4}; g) \ch{H2SeO3}; h) \ch{H3BO3}.
    \begin{solu}
        a) \textcolor{red}{dihidrogeno(trioxidocarbonato)}; b) \textcolor{red}{hidrogeno(trioxidoclorato)}; c)
        \textcolor{red}{hidrogeno(tetraoxidomanganato)}; d) \textcolor{red}{hidrogeno(oxidobromato)}; e)
        \textcolor{red}{dihidrogeno(tetraoxidosulfato)}; f) \textcolor{red}{tetrahidrogeno(tetraoxidosilicato)}; g)
        \textcolor{red}{dihidrogeno(trioxidoseleniato)}; h) \textcolor{red}{trihidrogeno(trioxidoborato)}.
    \end{solu}
\end{exercise}

\begin{exercise}
    \(\star\) Formula: a) hidrogeno(tetraoxidobromato); b) dihidrogeno(heptaoxidodisulfato); c) hidrogeno(oxidoclorato); d)
    trihidrogeno(trioxidoarseniato); e) dihidrogeno(tetraoxidoseleniato); f) hidrogeno(dioxidoyodato); g)
    tetrahidrogeno(heptaoxidodifosfato); h) dihidrogeno(trioxidotelurato).
    \begin{solu}
        a) \textcolor{red}{\ch{HBrO4}}; b) \textcolor{red}{\ch{H2S2O7}}; c) \textcolor{red}{\ch{HClO}}; d)
        \textcolor{red}{\ch{H3AsO3}}; e) \textcolor{red}{\ch{H2SeO4}}; f) \textcolor{red}{\ch{HIO2}}; g)
        \textcolor{red}{\ch{H4P2O7}}; h) \textcolor{red}{\ch{H2TeO3}}.
    \end{solu}
\end{exercise}

\begin{exercise}
    \(\star\star\) Nombra con la nomenclatura de adición y escribe el nombre tradicional: a) \ch{Br(OH)}; b) \ch{ClO2(OH)};
    c) \ch{MnO3(OH)}; d) \ch{IO(OH)}; e) \ch{NO(OH)}; f) \ch{CO(OH)2}; g) \ch{TeO2(OH)2}; h) \ch{As(OH)3}.
    \begin{solu}
        a) \textcolor{red}{hidroxidobromo; ácido hipobromoso} (\ch{HBrO}). b) \textcolor{red}{hidroxidodioxidocloro; ácido
        clórico} (\ch{HClO3}). c) \textcolor{red}{hidroxidotrioxidomanganeso; ácido permangánico} (\ch{HMnO4}). d)
        \textcolor{red}{hidroxidooxidoyodo; ácido yodoso} (\ch{HIO2}). e) \textcolor{red}{hidroxidooxidonitrógeno; ácido nitroso}
        (\ch{HNO2}). f) \textcolor{red}{dihidroxidooxidocarbono; ácido carbónico} (\ch{H2CO3}). g)
        \textcolor{red}{dihidroxidodioxidoteluro; ácido telúrico} (\ch{H2TeO4}). h) \textcolor{red}{trihidroxidoarsénico; ácido
        arsenoso} (\ch{H3AsO3}).
    \end{solu}
\end{exercise}

\begin{exercise}
    \(\star\star\) Formula: a) hidroxidotrioxidobromo; b) hidroxidooxidocloro; c) trihidroxidooxidofósforo; d)
    dihidroxidooxidoazufre; e) hidroxidodioxidonitrógeno; f) tetrahidroxidosilicio; g) dihidroxidodioxidocromo; h) hidroxidoyodo.
    \begin{solu}
        a) \textcolor{red}{\ch{HBrO4}}; b) \textcolor{red}{\ch{HClO2}}; c) \textcolor{red}{\ch{H3PO4}}; d)
        \textcolor{red}{\ch{H2SO3}}; e) \textcolor{red}{\ch{HNO3}}; f) \textcolor{red}{\ch{H4SiO4}}; g)
        \textcolor{red}{\ch{H2CrO4}}; h) \textcolor{red}{\ch{HIO}}.
    \end{solu}
\end{exercise}

\begin{exercise}
    \(\star\star\) Calcula el estado de oxidación del átomo central en los ácidos hipocloroso, cloroso, clórico y perclórico. ¿Qué
    relación hay entre los prefijos y sufijos del nombre tradicional y ese estado de oxidación?
    \begin{solu}
        \ch{HClO}: \textcolor{red}{\(+1\)}; \ch{HClO2}: \textcolor{red}{\(+3\)}; \ch{HClO3}: \textcolor{red}{\(+5\)}; \ch{HClO4}:
        \textcolor{red}{\(+7\)}. \textcolor{red}{De menor a mayor estado de oxidación: hipo-…-oso, -oso, -ico, per-…-ico}; cada
        oxígeno más sube el estado de oxidación del cloro en dos unidades.
    \end{solu}
\end{exercise}

\subsection{Oxoaniones}

\begin{exercise}
    \(\star\) Nombra con la nomenclatura de composición y da el nombre tradicional: a) \ch{NO3-}; b) \ch{ClO2-}; c) \ch{SO4^2-};
    d) \ch{PO3^3-}; e) \ch{IO3-}; f) \ch{BrO4-}; g) \ch{CrO4^2-}; h) \ch{AsO4^3-}.
    \begin{solu}
        a) \textcolor{red}{trioxidonitrato(1\nobreakdash--); nitrato}. b) \textcolor{red}{dioxidoclorato(1\nobreakdash--); clorito}. c)
        \textcolor{red}{tetraoxidosulfato(2\nobreakdash--); sulfato}. d) \textcolor{red}{trioxidofosfato(3\nobreakdash--); fosfito}. e)
        \textcolor{red}{trioxidoyodato(1\nobreakdash--); yodato}. f) \textcolor{red}{tetraoxidobromato(1\nobreakdash--); perbromato}. g)
        \textcolor{red}{tetraoxidocromato(2\nobreakdash--); cromato}. h) \textcolor{red}{tetraoxidoarseniato(3\nobreakdash--); arseniato}.
    \end{solu}
\end{exercise}

\begin{exercise}
    \(\star\) Formula: a) dioxidonitrato(1\nobreakdash--); b) trioxidobromato(1\nobreakdash--); c) trioxidoseleniato(2\nobreakdash--); d) tetraoxidosilicato(4\nobreakdash--);
    e) oxidoyodato(1\nobreakdash--); f) heptaoxidodicromato(2\nobreakdash--); g) hipoclorito; h) permanganato; i) carbonato; j) sulfito.
    \begin{solu}
        a) \textcolor{red}{\ch{NO2-}}; b) \textcolor{red}{\ch{BrO3-}}; c) \textcolor{red}{\ch{SeO3^2-}}; d)
        \textcolor{red}{\ch{SiO4^4-}}; e) \textcolor{red}{\ch{IO-}}; f) \textcolor{red}{\ch{Cr2O7^2-}}; g)
        \textcolor{red}{\ch{ClO-}}; h) \textcolor{red}{\ch{MnO4-}}; i) \textcolor{red}{\ch{CO3^2-}}; j) \textcolor{red}{\ch{SO3^2-}}.
    \end{solu}
\end{exercise}

\begin{exercise}
    \(\star\star\) Nombra con la nomenclatura de hidrógeno y con la de adición, y da el nombre tradicional: a) \ch{HSO3-}; b)
    \ch{H2BO3-}; c) \ch{HAsO4^2-}; d) \ch{HSeO4-}; e) \ch{H3SiO4-}.
    \begin{solu}
        a) \textcolor{red}{hidrogeno(trioxidosulfato)(1\nobreakdash--); \ch{SO2(OH)-}, hidroxidodioxidosulfato(1\nobreakdash--); hidrogenosulfito}.
        b) \textcolor{red}{dihidrogeno(trioxidoborato)(1\nobreakdash--); \ch{BO(OH)2-}, dihidroxidooxidoborato(1\nobreakdash--); dihidrogenoborato}.
        c) \textcolor{red}{hidrogeno(tetraoxidoarseniato)(2\nobreakdash--); \ch{AsO3(OH)^2-}, hidroxidotrioxidoarseniato(2\nobreakdash--);
        hidrogenoarseniato}. d) \textcolor{red}{hidrogeno(tetraoxidoseleniato)(1\nobreakdash--); \ch{SeO3(OH)-},
        hidroxidotrioxidoseleniato(1\nobreakdash--); hidrogenoseleniato}. e) \textcolor{red}{trihidrogeno(tetraoxidosilicato)(1\nobreakdash--);
        \ch{SiO(OH)3-}, trihidroxidooxidosilicato(1\nobreakdash--); trihidrogenosilicato}.
    \end{solu}
\end{exercise}

\begin{exercise}
    \(\star\star\) Formula: a) hidrogeno(trioxidocarbonato)(1\nobreakdash--); b) dihidrogeno(tetraoxidoarseniato)(1\nobreakdash--); c)
    hidroxidotrioxidofosfato(2\nobreakdash--); d) dihidroxidodioxidofosfato(1\nobreakdash--); e) hidroxidodioxidotelurato(1\nobreakdash--); f) hidrogenosulfato.
    \begin{solu}
        a) \textcolor{red}{\ch{HCO3-}}; b) \textcolor{red}{\ch{H2AsO4-}}; c) \textcolor{red}{\ch{HPO4^2-}}; d)
        \textcolor{red}{\ch{H2PO4-}}; e) \textcolor{red}{\ch{HTeO3-}}; f) \textcolor{red}{\ch{HSO4-}}.
    \end{solu}
\end{exercise}

\subsection{Oxosales}

\begin{exercise}
    \(\star\star\) Indica la carga del catión y del anión y el nombre tradicional de: a) \ch{KClO3}; b) \ch{Fe(NO2)2}; c)
    \ch{CuSO4}; d) \ch{Na3PO4}; e) \ch{Sn(SO4)2}; f) \ch{Ca(BrO)2}; g) \ch{Ag2CO3}; h) \ch{Al2(SiO3)3}.
    \begin{solu}
        a) \ch{K+} y \ch{ClO3-}: \textcolor{red}{clorato de potasio}. b) \ch{Fe^2+} y \ch{NO2-}: \textcolor{red}{nitrito de
        hierro(II)}. c) \ch{Cu^2+} y \ch{SO4^2-}: \textcolor{red}{sulfato de cobre(II)} (con \ch{Cu+} el S tendría \(+7\)). d)
        \ch{Na+} y \ch{PO4^3-}: \textcolor{red}{fosfato de sodio}. e) \ch{Sn^4+} y \ch{SO4^2-}: \textcolor{red}{sulfato de
        estaño(IV)}. f) \ch{Ca^2+} y \ch{BrO-}: \textcolor{red}{hipobromito de calcio}. g) \ch{Ag+} y \ch{CO3^2-}:
        \textcolor{red}{carbonato de plata}. h) \ch{Al^3+} y \ch{SiO3^2-}: \textcolor{red}{metasilicato de aluminio}
        (trioxidosilicato de aluminio).
    \end{solu}
\end{exercise}

\begin{exercise}
    \(\star\star\) Comprueba si son correctas estas oxosales, ayudándote de las cargas de los iones, y corrige las que no lo sean:
    a) carbonato de potasio, \ch{KCO3}; b) nitrato de cobre(II), \ch{Cu(NO3)2}; c) sulfato de hierro(III), \ch{Fe3(SO4)2}; d)
    permanganato de magnesio, \ch{MgMnO4}; e) fosfato de calcio, \ch{Ca3(PO4)2}; f) cromato de plomo(II), \ch{PbCrO4}.
    \begin{solu}
        a) \textcolor{red}{Incorrecta}: \ch{K+} y \ch{CO3^2-}, \textcolor{red}{\ch{K2CO3}}. b) \textcolor{red}{Correcta}. c)
        \textcolor{red}{Incorrecta}: \ch{Fe^3+} y \ch{SO4^2-}, \textcolor{red}{\ch{Fe2(SO4)3}}. d) \textcolor{red}{Incorrecta}:
        \ch{Mg^2+} y \ch{MnO4-}, \textcolor{red}{\ch{Mg(MnO4)2}}. e) \textcolor{red}{Correcta}. f) \textcolor{red}{Correcta}.
    \end{solu}
\end{exercise}

\begin{exercise}
    \(\star\star\) Nombra con la nomenclatura de composición y con la de adición: a) \ch{K2SO3}; b) \ch{Mg(NO3)2}; c)
    \ch{Fe2(CO3)3}; d) \ch{Ba3(PO4)2}; e) \ch{Cu(ClO4)2}; f) \ch{NaIO}.
    \begin{solu}
        a) \textcolor{red}{trioxidosulfato de dipotasio; trioxidosulfato(2\nobreakdash--) de potasio}. b) \textcolor{red}{bis(trioxidonitrato)
        de magnesio; trioxidonitrato(1\nobreakdash--) de magnesio}. c) \textcolor{red}{tris(trioxidocarbonato) de dihierro;
        trioxidocarbonato(2\nobreakdash--) de hierro(3+)}. d) \textcolor{red}{bis(tetraoxidofosfato) de tribario; tetraoxidofosfato(3\nobreakdash--) de
        bario}. e) \textcolor{red}{bis(tetraoxidoclorato) de cobre; tetraoxidoclorato(1\nobreakdash--) de cobre(2+)}. f)
        \textcolor{red}{oxidoyodato de sodio; oxidoyodato(1\nobreakdash--) de sodio}.
    \end{solu}
\end{exercise}

\begin{exercise}
    \(\star\star\) Formula: a) nitrato de calcio; b) sulfato de aluminio; c) clorato de sodio; d) carbonato de cobre(II); e)
    hipoclorito de potasio; f) dicromato de amonio; g) bis(trioxidobromato) de bario; h) tris(tetraoxidofosfato) de
    tetrahierro\dots{} ¿es posible?; i) tetraoxidosulfato(2\nobreakdash--) de cobalto(3+); j) trioxidonitrato(1\nobreakdash--) de plata; k) nitrito de
    hierro(III); l) yodato de mercurio(II).
    \begin{solu}
        a) \textcolor{red}{\ch{Ca(NO3)2}}; b) \textcolor{red}{\ch{Al2(SO4)3}}; c) \textcolor{red}{\ch{NaClO3}}; d)
        \textcolor{red}{\ch{CuCO3}}; e) \textcolor{red}{\ch{KClO}}; f) \textcolor{red}{\ch{(NH4)2Cr2O7}}; g)
        \textcolor{red}{\ch{Ba(BrO3)2}}; h) \textcolor{red}{\ch{Fe4(PO4)3}}: 9 cargas negativas entre 4 hierros no da un número
        entero (\(+9/4\)): \textcolor{red}{no es posible}. i) \textcolor{red}{\ch{Co2(SO4)3}}; j) \textcolor{red}{\ch{AgNO3}}; k)
        \textcolor{red}{\ch{Fe(NO2)3}}; l) \textcolor{red}{\ch{Hg(IO3)2}}.
    \end{solu}
\end{exercise}

\subsection{Sales ácidas y sales hidratadas}

\begin{exercise}
    \(\star\star\) Nombra (nombre tradicional y de composición): a) \ch{NaHCO3}; b) \ch{KHSO3}; c) \ch{Ca(H2PO4)2}; d)
    \ch{Fe(HSO4)2}; e) \ch{LiHS}; f) \ch{Al(HSe)3}; g) \ch{NH4HCO3}; h) \ch{CuHPO4}.
    \begin{solu}
        a) \textcolor{red}{hidrogenocarbonato de sodio; hidrogeno(trioxidocarbonato) de sodio}. b) \textcolor{red}{hidrogenosulfito
        de potasio; hidrogeno(trioxidosulfato) de potasio}. c) \textcolor{red}{dihidrogenofosfato de calcio;
        bis[dihidrogeno(tetraoxidofosfato)] de calcio}. d) \textcolor{red}{hidrogenosulfato de hierro(II);
        bis[hidrogeno(tetraoxidosulfato)] de hierro}. e) \textcolor{red}{hidrogenosulfuro de litio; hidrogeno(sulfuro) de litio}.
        f) \textcolor{red}{hidrogenoseleniuro de aluminio; tris[hidrogeno(seleniuro)] de aluminio}. g)
        \textcolor{red}{hidrogenocarbonato de amonio; hidrogeno(trioxidocarbonato) de amonio}. h) \textcolor{red}{hidrogenofosfato
        de cobre(II); hidrogeno(tetraoxidofosfato) de cobre}.
    \end{solu}
\end{exercise}

\begin{exercise}
    \(\star\star\) Formula: a) hidrogenosulfato de sodio; b) hidrogenocarbonato de magnesio; c) dihidrogenofosfato de amonio; d)
    hidrogenosulfuro de calcio; e) hidrogenofosfato de potasio; f) bis[hidrogeno(telururo)] de cinc; g) hidrogenoarseniato de
    hierro(II); h) hidrogenosulfito de plomo(IV).
    \begin{solu}
        a) \textcolor{red}{\ch{NaHSO4}}; b) \textcolor{red}{\ch{Mg(HCO3)2}}; c) \textcolor{red}{\ch{NH4H2PO4}}; d)
        \textcolor{red}{\ch{Ca(HS)2}}; e) \textcolor{red}{\ch{K2HPO4}}; f) \textcolor{red}{\ch{Zn(HTe)2}}; g)
        \textcolor{red}{\ch{FeHAsO4}}; h) \textcolor{red}{\ch{Pb(HSO3)4}}.
    \end{solu}
\end{exercise}

\begin{exercise}
    \(\star\) Nombra o formula estas sales hidratadas: a) \ch{CuSO4 * 5 H2O}; b) \ch{Na2CO3 * 10 H2O}; c) \ch{NiCl2 * 6 H2O}; d)
    \ch{2 CaSO4 * H2O}; e) sulfato de hierro(II) heptahidrato; f) cloruro de magnesio---agua (1/6); g) nitrato de cobre(II)
    trihidrato.
    \begin{solu}
        a) \textcolor{red}{sulfato de cobre(II) pentahidrato}. b) \textcolor{red}{carbonato de sodio decahidrato}. c)
        \textcolor{red}{cloruro de níquel(II) hexahidrato}. d) \textcolor{red}{sulfato de calcio---agua (2/1)}. e)
        \textcolor{red}{\ch{FeSO4 * 7 H2O}}; f) \textcolor{red}{\ch{MgCl2 * 6 H2O}}; g) \textcolor{red}{\ch{Cu(NO3)2 * 3 H2O}}.
    \end{solu}
\end{exercise}

\begin{exercise}
    \(\star\star\) Al calentar \qty{5.00}{\gram} de cloruro de bario hidratado hasta eliminar toda el agua quedan
    \qty{4.26}{\gram} de \ch{BaCl2}. Halla la fórmula de la sal hidratada y nómbrala (\(M(\ch{BaCl2}) = \qty{208.2}{\gram\per\mol}\)).
    \begin{solu}
        \(m(\ch{H2O}) = 5{,}00 - 4{,}26 = \qty{0.74}{\gram}\). \(n(\ch{BaCl2}) = \dfrac{4{,}26}{208{,}2} = \qty{0.0205}{\mol}\);
        \(n(\ch{H2O}) = \dfrac{0{,}74}{18{,}0} = \qty{0.0411}{\mol}\). Relación \(\dfrac{0{,}0411}{0{,}0205} = 2{,}0\):
        \textcolor{red}{\ch{BaCl2 * 2 H2O}, cloruro de bario dihidrato}.
    \end{solu}
\end{exercise}

%========================================================================================
\section{Ejercicios de ampliación}\label{sec:ud5-ampliacion}
%========================================================================================

\subsection{Compuestos binarios}

\begin{exercise}
    \(\star\star\) Nombra de todas las maneras posibles: a) \ch{CoO}; b) \ch{Au2O}; c) \ch{SrO}; d) \ch{O3Br2}; e) \ch{SeO3};
    f) \ch{K2O2}; g) \ch{BaH2}; h) \ch{CuH}; i) \ch{FeI3}; j) \ch{HBr}; k) \ch{CdO2}; l) \ch{PbCl2}; m) \ch{Mg2Si}; n)
    \ch{CS2}; ñ) \ch{NiS}; o) \ch{SbH3}; p) \ch{N2O}; q) \ch{Ga2S3}.
    \begin{solu}
        a) \textcolor{red}{monóxido de cobalto; óxido de cobalto(II); óxido de cobalto(2+)}. b) \textcolor{red}{óxido de dioro;
        óxido de oro(I); óxido de oro(1+)}. c) \textcolor{red}{óxido de estroncio}. d) \textcolor{red}{dibromuro de trioxígeno}.
        e) \textcolor{red}{trióxido de selenio; óxido de selenio(VI)}. f) \textcolor{red}{dióxido de dipotasio; peróxido de
        potasio}. g) \textcolor{red}{dihidruro de bario; hidruro de bario}. h) \textcolor{red}{hidruro de cobre(I); hidruro de
        cobre(1+)}. i) \textcolor{red}{triyoduro de hierro; yoduro de hierro(III); yoduro de hierro(3+)}. j)
        \textcolor{red}{bromuro de hidrógeno; bromano}. k) \textcolor{red}{dióxido de cadmio; peróxido de cadmio}. l)
        \textcolor{red}{dicloruro de plomo; cloruro de plomo(II); cloruro de plomo(2+)}. m) \textcolor{red}{siliciuro de
        dimagnesio; siliciuro de magnesio}. n) \textcolor{red}{disulfuro de carbono; sulfuro de carbono(IV)}. ñ)
        \textcolor{red}{sulfuro de níquel; sulfuro de níquel(II); sulfuro de níquel(2+)}. o) \textcolor{red}{trihidruro de
        antimonio; hidruro de antimonio(III); estibano}. p) \textcolor{red}{monóxido de dinitrógeno; óxido de nitrógeno(I)}. q)
        \textcolor{red}{trisulfuro de digalio; sulfuro de galio}.
    \end{solu}
\end{exercise}

\begin{exercise}
    \(\star\star\) Formula: a) óxido de estroncio; b) óxido de mercurio(I); c) óxido de fósforo(V); d) dicloruro de
    pentaoxígeno; e) dióxido de dicesio; f) hidruro de estaño(II); g) seleniuro de dihidrógeno; h) amoniaco; i) fosfuro de
    potasio; j) peróxido de rubidio; k) peróxido de plata; l) óxido de níquel(2+); m) fluoruro de cobre(II); n) tetracloruro de
    titanio; ñ) nitruro de magnesio; o) sulfuro de cobalto(III); p) hidruro de titanio(III); q) telururo de cadmio; r) carburo
    de silicio.
    \begin{solu}
        a) \textcolor{red}{\ch{SrO}}; b) \textcolor{red}{\ch{Hg2O}}; c) \textcolor{red}{\ch{P2O5}}; d) \textcolor{red}{\ch{O5Cl2}};
        e) \textcolor{red}{\ch{Cs2O2}}; f) \textcolor{red}{\ch{SnH2}}; g) \textcolor{red}{\ch{H2Se}}; h) \textcolor{red}{\ch{NH3}};
        i) \textcolor{red}{\ch{K3P}}; j) \textcolor{red}{\ch{Rb2O2}}; k) \textcolor{red}{\ch{Ag2O2}}; l) \textcolor{red}{\ch{NiO}};
        m) \textcolor{red}{\ch{CuF2}}; n) \textcolor{red}{\ch{TiCl4}}; ñ) \textcolor{red}{\ch{Mg3N2}}; o)
        \textcolor{red}{\ch{Co2S3}}; p) \textcolor{red}{\ch{TiH3}}; q) \textcolor{red}{\ch{CdTe}}; r) \textcolor{red}{\ch{SiC}}.
    \end{solu}
\end{exercise}

\subsection{Compuestos ternarios y cuaternarios}

\begin{exercise}
    \(\star\star\star\) Nombra de todas las maneras posibles que conozcas: a) \ch{Co(OH)2}; b) \ch{CuOH}; c) \ch{H2CO3}; d)
    \ch{HClO2}; e) \ch{H2SeO4}; f) \ch{HNO2}; g) \ch{IO4-}; h) \ch{HPO4^2-}; i) \ch{NaNO3}; j) \ch{Ni3(PO4)2}; k)
    \ch{Hg(ClO)2}; l) \ch{KHSO4}; m) \ch{Cu(HCO3)2}; n) \ch{Fe(HS)3}; ñ) \ch{CaCrO4}; o) \ch{K2MnO4}.
    \begin{solu}
        a) \textcolor{red}{dihidróxido de cobalto; hidróxido de cobalto(II); hidróxido de cobalto(2+)}. b) \textcolor{red}{hidróxido
        de cobre(I); hidróxido de cobre(1+)}. c) \textcolor{red}{ácido carbónico; dihidrogeno(trioxidocarbonato);
        dihidroxidooxidocarbono}. d) \textcolor{red}{ácido cloroso; hidrogeno(dioxidoclorato); hidroxidooxidocloro}. e)
        \textcolor{red}{ácido selénico; dihidrogeno(tetraoxidoseleniato); dihidroxidodioxidoselenio}. f) \textcolor{red}{ácido
        nitroso; hidrogeno(dioxidonitrato); hidroxidooxidonitrógeno}. g) \textcolor{red}{peryodato; tetraoxidoyodato(1\nobreakdash--)}. h)
        \textcolor{red}{hidrogenofosfato; hidrogeno(tetraoxidofosfato)(2\nobreakdash--); hidroxidotrioxidofosfato(2\nobreakdash--)}. i)
        \textcolor{red}{nitrato de sodio; trioxidonitrato de sodio; trioxidonitrato(1\nobreakdash--) de sodio}. j) \textcolor{red}{fosfato de
        níquel(II); bis(tetraoxidofosfato) de triníquel; tetraoxidofosfato(3\nobreakdash--) de níquel(2+)}. k) \textcolor{red}{hipoclorito de
        mercurio(II); bis(oxidoclorato) de mercurio; oxidoclorato(1\nobreakdash--) de mercurio(2+)}. l) \textcolor{red}{hidrogenosulfato de
        potasio; hidrogeno(tetraoxidosulfato) de potasio}. m) \textcolor{red}{hidrogenocarbonato de cobre(II);
        bis[hidrogeno(trioxidocarbonato)] de cobre}. n) \textcolor{red}{hidrogenosulfuro de hierro(III); tris[hidrogeno(sulfuro)] de
        hierro}. ñ) \textcolor{red}{cromato de calcio; tetraoxidocromato de calcio; tetraoxidocromato(2\nobreakdash--) de calcio}. o)
        \textcolor{red}{manganato de potasio; tetraoxidomanganato de dipotasio; tetraoxidomanganato(2\nobreakdash--) de potasio}.
    \end{solu}
\end{exercise}

\begin{exercise}
    \(\star\star\star\) Formula: a) hidróxido de bario; b) hidróxido de hierro(3+); c) ácido sulfuroso; d) ácido perclórico; e)
    trihidrogeno(tetraoxidoarseniato); f) hidroxidodioxidobromo; g) dihidroxidooxidocarbono; h) hipobromito; i)
    tetraoxidofosfato(3\nobreakdash--); j) hidrogeno(trioxidosulfato)(1\nobreakdash--); k) dicromato de potasio; l) sulfato de amonio; m)
    tris(trioxidonitrato) de cromo; n) tetraoxidoseleniato(2\nobreakdash--) de cobre(1+); ñ) hidrogenocarbonato de calcio; o)
    dihidrogenofosfato de sodio; p) hidrogenoseleniuro de níquel(II); q) sulfato de cinc heptahidrato.
    \begin{solu}
        a) \textcolor{red}{\ch{Ba(OH)2}}; b) \textcolor{red}{\ch{Fe(OH)3}}; c) \textcolor{red}{\ch{H2SO3}}; d)
        \textcolor{red}{\ch{HClO4}}; e) \textcolor{red}{\ch{H3AsO4}}; f) \textcolor{red}{\ch{HBrO3}}; g) \textcolor{red}{\ch{H2CO3}};
        h) \textcolor{red}{\ch{BrO-}}; i) \textcolor{red}{\ch{PO4^3-}}; j) \textcolor{red}{\ch{HSO3-}}; k)
        \textcolor{red}{\ch{K2Cr2O7}}; l) \textcolor{red}{\ch{(NH4)2SO4}}; m) \textcolor{red}{\ch{Cr(NO3)3}}; n)
        \textcolor{red}{\ch{Cu2SeO4}}; ñ) \textcolor{red}{\ch{Ca(HCO3)2}}; o) \textcolor{red}{\ch{NaH2PO4}}; p)
        \textcolor{red}{\ch{Ni(HSe)2}}; q) \textcolor{red}{\ch{ZnSO4 * 7 H2O}}.
    \end{solu}
\end{exercise}

\begin{exercise}
    \(\star\star\star\) Un compuesto tiene la fórmula \ch{MXO3}, donde M es un metal alcalino. a) ¿Qué estado de oxidación tiene
    X? b) Propón un elemento X posible y nombra el compuesto. c) Si M fuera un metal alcalinotérreo, ¿cómo cambiaría la fórmula
    para conservar el mismo anión?
    \begin{solu}
        a) \(+1 + x - 6 = 0\): \textcolor{red}{\(x = +5\)}. b) \textcolor{red}{X puede ser N, Cl, Br o I}: por ejemplo,
        \ch{KNO3}, nitrato de potasio, o \ch{NaClO3}, clorato de sodio. c) El anión \ch{XO3-} tiene carga \(-1\) y el catión
        \(+2\): \textcolor{red}{\ch{M(XO3)2}}, por ejemplo \ch{Ca(NO3)2}.
    \end{solu}
\end{exercise}

\begin{exercise}
    \(\star\star\star\) El yeso (\ch{CaSO4 * 2 H2O}) se transforma al calentarlo en escayola, \ch{2 CaSO4 * H2O}. a) Nombra
    las dos sustancias. b) ¿Cuántos gramos de agua pierde 1 kg de yeso al convertirse en escayola? (\(M(\ch{CaSO4}) =
    \qty{136.2}{\gram\per\mol}\)).
    \begin{solu}
        a) \textcolor{red}{Sulfato de calcio dihidrato y sulfato de calcio---agua (2/1)} (a veces llamado hemihidrato).
        b) Por cada 2 \ch{CaSO4} se pasa de 4 \ch{H2O} a 1: se pierden 1,5 \ch{H2O} por \ch{CaSO4}.
        \(n = \dfrac{1000}{172{,}2} = \qty{5.81}{\mol}\) de yeso; \(m = 5{,}81 \cdot 1{,}5 \cdot 18{,}0 =
        \textcolor{red}{\qty{157}{\gram}}\) de agua.
    \end{solu}
\end{exercise}

\begin{exercise}
    \(\star\star\star\) Busca el error en cada fórmula o nombre y corrígelo: a) \ch{NaOH2}, hidróxido de sodio; b) \ch{H2NO3},
    ácido nítrico; c) \ch{Ca(CO3)2}, carbonato de calcio; d) \ch{Fe2O3}, trióxido de hierro; e) \ch{HSO4}, hidrogenosulfato; f)
    \ch{Cl2O}, monóxido de dicloro.
    \begin{solu}
        a) \textcolor{red}{\ch{NaOH}}: el sodio tiene \(+1\) y basta un \ch{OH-}. b) \textcolor{red}{\ch{HNO3}}: con 2 H, el N
        tendría \(+4\). c) \textcolor{red}{\ch{CaCO3}}: \ch{Ca^2+} y \ch{CO3^2-}, uno de cada. d) \textcolor{red}{Trióxido de
        dihierro}: hay que indicar también los 2 hierros. e) \textcolor{red}{\ch{HSO4-}}: le falta la carga. f)
        \textcolor{red}{\ch{OCl2}, dicloruro de oxígeno}: el O va antes que el Cl en la secuencia y el compuesto es un haluro de
        oxígeno.
    \end{solu}
\end{exercise}

\fussy
